\documentclass[11pt,a4paper]{article}
\usepackage[T1]{fontenc}
\usepackage{lmodern,amsmath,amssymb,amsthm,mathtools,microtype}
\usepackage[margin=27mm]{geometry}
\usepackage[colorlinks=true,linkcolor=blue,citecolor=blue,urlcolor=blue]{hyperref}
\setlength{\emergencystretch}{2em}
\newtheorem{proposition}{Proposition}[section]
\newtheorem{example}[proposition]{Example}
\newcommand{\Q}{\mathbb Q}\newcommand{\Z}{\mathbb Z}
\DeclareMathOperator{\Tr}{Tr}\DeclareMathOperator{\Nm}{N}
\title{Rational points in a rescaled native square incidence}
\author{}\date{17 November 2026}
\begin{document}\maketitle
\begin{abstract}
A tangent construction gives an explicit rational point on the native
square incidence when the physical weight can be rescaled. Its positive
branch produces an exact prescribed-trace example on the elliptic
multiplication family at inputs $\pm2$. The rescaling is a fixed rational
constant for this example. Uniform realization with finitely many fixed
rescalings remains an additional arithmetic requirement. We also give
the exact genus-three equation obtained when the positive prefactor is
a rational square.
\end{abstract}

\section{A positive tangent section}

Let $W=a+b\sqrt2\in\Q(\sqrt2)$, and put
\[
 N=a^2-2b^2,\qquad H=N-a,\qquad J=N-2a+1.
\]
For a rational physical weight $w$ write $g=1+w$ and
\[
 \Delta=g^2N-2ga+1,\qquad
 {\cal S}=\frac{w(gN-a)}{\Delta}.
\]
On the guarded domain the native radical of \cite{Single} satisfies
\begin{equation}\label{eq:equiv}
 {\cal R}={\cal S}\left(\frac{\Delta}{wgN}\right)^2.
\end{equation}

\begin{proposition}\label{prop:tangent}
Suppose that both real values of $W$ exceed one and
$8b^2>3H^2$. Then
\[
 w_* =\frac{4HJ}{8b^2-3H^2}>0,\qquad
 k_* =\frac{2H^2}{8b^2-H^2}\in(0,1)
\]
satisfy ${\cal S}(w_*,W)=k_*^2$. Consequently
\[
 \frac{k_*\Delta(w_*)}{w_*(1+w_*)N}
\]
is a nonzero rational square root of the full native radical.
\end{proposition}
\begin{proof}
The equation ${\cal S}=k^2$ is quadratic in $w$:
\[
 N(1-k^2)w^2+H(1-2k^2)w-Jk^2=0.
\]
Using $H^2-NJ=2b^2$, its discriminant is
\[
 y^2=H^2+8b^2 k^2(k^2-1).
\]
The parabola
\[
 y=H+\frac{8b^2}{H}(k^2-k)
\]
passes through $(0,H)$ and is tangent at $(1,H)$. Substitution gives
\[
 y^2-H^2-8b^2k^2(k^2-1)
 =\frac{8b^2}{H^2}k(k-1)^2\bigl((8b^2-H^2)k-2H^2\bigr).
\]
The fourth intersection yields $k_*$, and the positive quadratic
branch yields $w_*$. Direct substitution verifies ${\cal S}=k_*^2$.

Write the two values of $W$ as $W_+,W_->1$. Then
$J=(W_+-1)(W_--1)>0$ and
$H=J+(W_++W_-)/2-1>0$. Moreover
$8b^2=(W_+-W_-)^2<4H^2$. Together with the assumed lower bound,
this gives $0<k_*<1$. Positivity of $w_*$ follows directly.
Each $gW_\pm-1$ is positive, so $\Delta>0$. Equation
\eqref{eq:equiv} gives the final assertion.
\end{proof}

\section{An actual multiplication-family example}

Use the curve and source functions of \cite{Single}:
\[
 E:y^2=x^3-(3-\sqrt2)^2x,\qquad
 P=\left(\frac{57-28\sqrt2}{8},\frac{-660+417\sqrt2}{32}\right).
\]
Put $f=1/x$, $\zeta=-x/y$, and retain the six tests
\[
 \begin{gathered}
 r=\zeta(U)/\zeta(H_1),\qquad
 s=\zeta(U-A)/\zeta(H_1-P),\\
 C_1=f(H_1),\quad C_2=f(U),\quad
 C_3=\overline{f(H_1-P)},\quad C_4=\overline{f(U-A)},\\
 C_5=r-\bar s,\qquad C_6=r+\bar s-2X.
 \end{gathered}
\]
The notation $H_1$ distinguishes the elliptic reference point from
$H=N-a$ in Proposition~\ref{prop:tangent}.

\begin{example}\label{ex:two}
Take
\[
 X=2,\quad A=2P,\quad H_1=4P,\quad U=8P=[X]H_1,
\]
and use the ordinary-$abc$ height parameters
\[
 B=3,\quad e=8,\quad d=6,\quad z=1,\qquad
 t=\left(7^8+(\Tr x(2P))^8\right)^6+1.
\]
Define the actual source values
\[
 F(t)=\frac{t(t+1)(2t+1)(3t+1)(4t+1)(5t+2)}{120},
 \quad w_0=\frac{F(t)}{4(t-1)^8},\quad
 W=\prod_{i=1}^6\left(1+\frac{C_i^2}{F(t)}\right).
\]
The exact certificate accompanying this manuscript verifies
$8b^2>3H^2$ for $W=a+b\sqrt2$. Fix the rational constant
\begin{equation}\label{eq:gamma}
 \gamma=\frac{w_*}{w_0}
 =\frac{4HJ}{(8b^2-3H^2)w_0}.
\end{equation}
Then the actual radical with physical weight $\gamma w_0$ is a
nonzero rational square, with the square root in
Proposition~\ref{prop:tangent}.

Every retained point guard is satisfied. The differences at all
ordinate-sign characters have nonzero labels $3,5,6,10$, and the
three original point charts are checked exactly. At every character
the detector slots exceed one. The fixed positive weight therefore
makes both native radical factors positive and nonzero. These facts
include all final character guards.

Replacing $(X,A,U)$ by $(-X,-A,-U)$ changes the signs of $C_5,C_6$
and preserves every other squared test, the trace, and the argument.
Thus the same frozen constant also gives an example at $X=-2$.
\end{example}

The certificate uses exact rational arithmetic and the complete
source equations. Its reduced $t$ has numerator and denominator
lengths $478$ and $435$ decimal digits. The corresponding lengths
for $\gamma$ are $70136$ and $70308$. Equation~\eqref{eq:gamma},
the recorded rational tests, and the checker give an exact finite
specification of this constant.

\section{Fixed rescalings and the remaining arithmetic condition}

For a fixed exceptional set $S$, a fixed rational rescaling $\gamma$
preserves every consumed valuation if $v_p(\gamma)=0$ for every
$p\notin S$. If this condition is not known, one may instead fix
\[
 S^*=S\cup\{p:v_p(\gamma)\ne0\}
\]
once and for all. The contact-height bound changes by a fixed additive
constant, and every local extraction outside $S^*$ is unchanged.
The coefficient $\gamma$ may not then vary with the input.

More generally, a finite menu $\Gamma\subset\Q_{>0}$ is compatible
with the same argument after its finite prime supports are included
in a common fixed exceptional set. Uniform completeness would require
that every integer input have an actual guarded source tuple for which
the full native radical is a square for some $\gamma\in\Gamma$.
The tangent section gives the exact sufficient condition
\[
 \Gamma\cap
 \left\{
 \frac{4H(W)J(W)}{(8b(W)^2-3H(W)^2)w_0}:
 \text{actual guarded tuples at }X,\quad8b(W)^2>3H(W)^2
 \right\}\ne\varnothing
\]
for every such input. The example proves this intersection for
$X=\pm2$ with a one-element menu. No fixed finite menu with the
required intersections for all integer inputs has been constructed.
In particular, this example is not a defining polynomial with a
smaller variable census. The shift $B=3$ is an algebraic example,
not a certified sufficient defining shift.

\section{A square positive prefactor}

The alternative normalizer
\[
 w=\frac{4u}{(1-u)^2},\qquad
 g=1+w=\left(\frac{1+u}{1-u}\right)^2
\]
preserves every consumed odd valuation outside two: if $v_p(u)$ is
odd then $v_p(w)=|v_p(u)|$. Values introduced by $1-u$ have even
valuation. The guards $u\ne1,-1$ hold in a prepared neighborhood
of $u=0$. With the fast conditional detector, its product remains
unchanged, because it uses $1/F(t)$ rather than $w$.

Write $g=s_0^2$. Clearing a rational square factor in the native
condition gives the exact equation
\begin{equation}\label{eq:genus3}
 Y^2=(s_0^2-1)(Ns_0^2-a)(Ns_0^4-2as_0^2+1).
\end{equation}
The discriminant of its degree-eight right side is
\[
 4096\,aN^6(a-N)^4(a^2-N)^6(2a-N-1)^4.
\]
For distinct slots $W_\pm>1$, every displayed factor is nonzero;
thus \eqref{eq:genus3} is a smooth genus-three curve. Its quotient
under $s_0\mapsto-s_0$ is the original genus-one weight incidence.

At $X=\pm1$, take $A=\pm P$ and $U=\pm H_1$. The comparison
tests vanish and
\[
 W=Z^2,\qquad
 Z=\left(1+\frac{f(H_1)^2}{F(t)}\right)
   \left(1+\frac{\overline{f(H_1-P)}^2}{F(t)}\right).
\]
Writing $Z=c+d_1\sqrt2$ gives $N=(c^2-2d_1^2)^2$ and
$a=c^2+2d_1^2$. The final quartic factor in
\eqref{eq:genus3} splits as
$(s_0Z-1)(s_0Z+1)(s_0\bar Z-1)(s_0\bar Z+1)$ over
$\Q(\sqrt2)$. This factorization does not remove the prescribed
square-abscissa lifting condition. The evident points $s_0=\pm1$
and the two points at infinity lie on guarded boundary loci.
No finite rational witness for the prescribed $s_0$ follows from
the simultaneous square structures alone.

\begin{thebibliography}{9}\raggedright
\bibitem{Single}\emph{A single-cover criterion from a mixed-power height
inequality}, repository manuscript,
\nolinkurl{papers/pole_parameters/single_cover_criterion_20261117_v56.tex}.
\end{thebibliography}
\end{document}
